\lecture{8}
\title{Systems}

\begin{document}

\subsection{Lecture and Recitation: Systems}

\textbf{Definition.} A \textbf{system} is a function
that maps signals $x[n]$ to signals $y[n]$.
\begin{itemize}
    \item Such a system is \textbf{additive} if
    $x_1[n] + x_2[n] \mapsto y_1[n] + y_2[n]$.
    \item Such a system  is \textbf{homogeneous} if
    $\lambda x[n] \mapsto \lambda y[n]$ for any constant $\lambda$.
    \item Such a system  is \textbf{linear} if
    it is both additive and homogeneous.
    \item Such a system  is \textbf{time-invariant} if
    $x[n - n_0] \mapsto y[n - n_0]$ for all integers $n_0$.
\end{itemize}

\begin{remark}
    We call the output signal $y[n]$ of a system
    the \textbf{response} to the input signal $x[n]$.

    One might imagine the input signal as, say,
    the voltage supplied to an AC circuit,
    and the output signal as, say,
    the voltage across a resistor.
    In this sense,
    the output signal is a \emph{response}
    to the input signal.
\end{remark}

\textbf{Example.} Consider the following systems.
\begin{itemize}
    \item The system $x[n] \mapsto x[n - 1] + x[n] + x[n + 1]$
    is linear and time-invariant.
    \item The system $x[n] \mapsto nx[n]$
    is linear but not time-invariant.
    \item The system $x[n] \mapsto x[n] + 1$
    is not linear but time-invariant.
    \item The system $x[n] \mapsto 0$ (the constant $0$ signal)
    is linear and time-invariant.
\end{itemize}

\textbf{Definition.} A system is said to be \textbf{LTI}
if it is both linear (L) and time-invariant (TI).

From now on, we will restrict our attention to LTI systems.

\end{document}